% Options for packages loaded elsewhere
% Options for packages loaded elsewhere
\PassOptionsToPackage{unicode}{hyperref}
\PassOptionsToPackage{hyphens}{url}
\PassOptionsToPackage{dvipsnames,svgnames,x11names}{xcolor}
%
\documentclass[
  a4paper,
  DIV=11,
  numbers=noendperiod,
  oneside]{scrreport}
\usepackage{xcolor}
\usepackage[left=10mm,marginparwidth=60mm,textwidth=130mm,marginparsep=5mm]{geometry}
\usepackage{amsmath,amssymb}
\setcounter{secnumdepth}{-\maxdimen} % remove section numbering
\usepackage{iftex}
\ifPDFTeX
  \usepackage[T1]{fontenc}
  \usepackage[utf8]{inputenc}
  \usepackage{textcomp} % provide euro and other symbols
\else % if luatex or xetex
  \usepackage{unicode-math} % this also loads fontspec
  \defaultfontfeatures{Scale=MatchLowercase}
  \defaultfontfeatures[\rmfamily]{Ligatures=TeX,Scale=1}
\fi
\usepackage{lmodern}
\ifPDFTeX\else
  % xetex/luatex font selection
\fi
% Use upquote if available, for straight quotes in verbatim environments
\IfFileExists{upquote.sty}{\usepackage{upquote}}{}
\IfFileExists{microtype.sty}{% use microtype if available
  \usepackage[]{microtype}
  \UseMicrotypeSet[protrusion]{basicmath} % disable protrusion for tt fonts
}{}
\makeatletter
\@ifundefined{KOMAClassName}{% if non-KOMA class
  \IfFileExists{parskip.sty}{%
    \usepackage{parskip}
  }{% else
    \setlength{\parindent}{0pt}
    \setlength{\parskip}{6pt plus 2pt minus 1pt}}
}{% if KOMA class
  \KOMAoptions{parskip=half}}
\makeatother
% Make \paragraph and \subparagraph free-standing
\makeatletter
\ifx\paragraph\undefined\else
  \let\oldparagraph\paragraph
  \renewcommand{\paragraph}{
    \@ifstar
      \xxxParagraphStar
      \xxxParagraphNoStar
  }
  \newcommand{\xxxParagraphStar}[1]{\oldparagraph*{#1}\mbox{}}
  \newcommand{\xxxParagraphNoStar}[1]{\oldparagraph{#1}\mbox{}}
\fi
\ifx\subparagraph\undefined\else
  \let\oldsubparagraph\subparagraph
  \renewcommand{\subparagraph}{
    \@ifstar
      \xxxSubParagraphStar
      \xxxSubParagraphNoStar
  }
  \newcommand{\xxxSubParagraphStar}[1]{\oldsubparagraph*{#1}\mbox{}}
  \newcommand{\xxxSubParagraphNoStar}[1]{\oldsubparagraph{#1}\mbox{}}
\fi
\makeatother

\usepackage{color}
\usepackage{fancyvrb}
\newcommand{\VerbBar}{|}
\newcommand{\VERB}{\Verb[commandchars=\\\{\}]}
\DefineVerbatimEnvironment{Highlighting}{Verbatim}{commandchars=\\\{\}}
% Add ',fontsize=\small' for more characters per line
\usepackage{framed}
\definecolor{shadecolor}{RGB}{241,243,245}
\newenvironment{Shaded}{\begin{snugshade}}{\end{snugshade}}
\newcommand{\AlertTok}[1]{\textcolor[rgb]{0.68,0.00,0.00}{#1}}
\newcommand{\AnnotationTok}[1]{\textcolor[rgb]{0.37,0.37,0.37}{#1}}
\newcommand{\AttributeTok}[1]{\textcolor[rgb]{0.40,0.45,0.13}{#1}}
\newcommand{\BaseNTok}[1]{\textcolor[rgb]{0.68,0.00,0.00}{#1}}
\newcommand{\BuiltInTok}[1]{\textcolor[rgb]{0.00,0.23,0.31}{#1}}
\newcommand{\CharTok}[1]{\textcolor[rgb]{0.13,0.47,0.30}{#1}}
\newcommand{\CommentTok}[1]{\textcolor[rgb]{0.37,0.37,0.37}{#1}}
\newcommand{\CommentVarTok}[1]{\textcolor[rgb]{0.37,0.37,0.37}{\textit{#1}}}
\newcommand{\ConstantTok}[1]{\textcolor[rgb]{0.56,0.35,0.01}{#1}}
\newcommand{\ControlFlowTok}[1]{\textcolor[rgb]{0.00,0.23,0.31}{\textbf{#1}}}
\newcommand{\DataTypeTok}[1]{\textcolor[rgb]{0.68,0.00,0.00}{#1}}
\newcommand{\DecValTok}[1]{\textcolor[rgb]{0.68,0.00,0.00}{#1}}
\newcommand{\DocumentationTok}[1]{\textcolor[rgb]{0.37,0.37,0.37}{\textit{#1}}}
\newcommand{\ErrorTok}[1]{\textcolor[rgb]{0.68,0.00,0.00}{#1}}
\newcommand{\ExtensionTok}[1]{\textcolor[rgb]{0.00,0.23,0.31}{#1}}
\newcommand{\FloatTok}[1]{\textcolor[rgb]{0.68,0.00,0.00}{#1}}
\newcommand{\FunctionTok}[1]{\textcolor[rgb]{0.28,0.35,0.67}{#1}}
\newcommand{\ImportTok}[1]{\textcolor[rgb]{0.00,0.46,0.62}{#1}}
\newcommand{\InformationTok}[1]{\textcolor[rgb]{0.37,0.37,0.37}{#1}}
\newcommand{\KeywordTok}[1]{\textcolor[rgb]{0.00,0.23,0.31}{\textbf{#1}}}
\newcommand{\NormalTok}[1]{\textcolor[rgb]{0.00,0.23,0.31}{#1}}
\newcommand{\OperatorTok}[1]{\textcolor[rgb]{0.37,0.37,0.37}{#1}}
\newcommand{\OtherTok}[1]{\textcolor[rgb]{0.00,0.23,0.31}{#1}}
\newcommand{\PreprocessorTok}[1]{\textcolor[rgb]{0.68,0.00,0.00}{#1}}
\newcommand{\RegionMarkerTok}[1]{\textcolor[rgb]{0.00,0.23,0.31}{#1}}
\newcommand{\SpecialCharTok}[1]{\textcolor[rgb]{0.37,0.37,0.37}{#1}}
\newcommand{\SpecialStringTok}[1]{\textcolor[rgb]{0.13,0.47,0.30}{#1}}
\newcommand{\StringTok}[1]{\textcolor[rgb]{0.13,0.47,0.30}{#1}}
\newcommand{\VariableTok}[1]{\textcolor[rgb]{0.07,0.07,0.07}{#1}}
\newcommand{\VerbatimStringTok}[1]{\textcolor[rgb]{0.13,0.47,0.30}{#1}}
\newcommand{\WarningTok}[1]{\textcolor[rgb]{0.37,0.37,0.37}{\textit{#1}}}

\usepackage{longtable,booktabs,array}
\usepackage{calc} % for calculating minipage widths
% Correct order of tables after \paragraph or \subparagraph
\usepackage{etoolbox}
\makeatletter
\patchcmd\longtable{\par}{\if@noskipsec\mbox{}\fi\par}{}{}
\makeatother
% Allow footnotes in longtable head/foot
\IfFileExists{footnotehyper.sty}{\usepackage{footnotehyper}}{\usepackage{footnote}}
\makesavenoteenv{longtable}
\usepackage{graphicx}
\makeatletter
\newsavebox\pandoc@box
\newcommand*\pandocbounded[1]{% scales image to fit in text height/width
  \sbox\pandoc@box{#1}%
  \Gscale@div\@tempa{\textheight}{\dimexpr\ht\pandoc@box+\dp\pandoc@box\relax}%
  \Gscale@div\@tempb{\linewidth}{\wd\pandoc@box}%
  \ifdim\@tempb\p@<\@tempa\p@\let\@tempa\@tempb\fi% select the smaller of both
  \ifdim\@tempa\p@<\p@\scalebox{\@tempa}{\usebox\pandoc@box}%
  \else\usebox{\pandoc@box}%
  \fi%
}
% Set default figure placement to htbp
\def\fps@figure{htbp}
\makeatother





\setlength{\emergencystretch}{3em} % prevent overfull lines

\providecommand{\tightlist}{%
  \setlength{\itemsep}{0pt}\setlength{\parskip}{0pt}}



 


% This is used for pdf compilation from qmd.
% Any special LaTeX commands should be updated in parallel with mymacros.sty which is used for html compilation.
% All rows *might* need an HTML equivalent. Those at the bottom of this file are identical.
% This is how the two files are included in the _quarto.yaml header of the qmd file
%format:
%  html:
%    defaults: [resources/latex/global-maths-definitions.yaml]
%  pdf:
%   include-in-header: ["include-in-header.tex"]

% the 'defaults' above secretly points to mymacros.sty where they really live.

\usepackage{gensymb}
\usepackage{amsmath}
\usepackage{xcolor}

\usepackage[skins,breakable]{tcolorbox}

% Code to reduce width of code blocks in PDF to encourage line wrapping
\usepackage{fvextra}
      \DefineVerbatimEnvironment{Highlighting}{Verbatim}{
        breaklines,
        breakanywhere,
        commandchars=\\\{\},
        fontsize=\small
      }
% Also line wrap in verbatim blocks
\RecustomVerbatimEnvironment{verbatim}{Verbatim}{breaklines,fontsize=\footnotesize}

\usepackage{environ}

% Below here is an old way to defining new custom coloured boxes.
% We use nice extensions now.

% \definecolor{goals-col}{RGB}{201, 232, 222}
% 
% % For defining new callout boxes
% % HTML version is in CSS file
% % Need to turn on goals and goals-header as environments using the latex-environment extension for quarto
% % This is designed to work for the following fenced div syntax:
% % :::{.goals}
% % ::::{.goals-header}
% % TITLE
% % ::::
% % Actual box contents
% % :::
% \newtcolorbox{goals}[1][]{notitle, top=0pt, boxsep=0pt, coltitle=black!90!white, colback=white, colframe=goals-col,#1}
% \NewEnviron{goals-header}{\tcbsubtitle[top=1mm, bottom=1mm]{\BODY}}
% % It works by creating a box with no title. Then no spacing. Then a subtitle defined later.

% Increasing spacing between rows of tables
\def\arraystretch{2}

% Syntax for highlighting my KEY TERMS
% HTML version is included in css for .term
\newcommand{\term}[1]{\textbf{#1}}

% Added to do derivatives, custom commands for HTML
% read by MathJax are inserted via diff.macro
% must keep naming in sync with esdiff package
\usepackage{esdiff}

%%%%%%%%%%%%%%%%%%%%%%%%%%%%%%%%%%%%%%%%%%%%%%%%%%%%%%%%%%%%
%%% BELOW HERE ARE COMMANDS SHARED WITH THE HTML OUTPUT %%%%
%%%%%%%%%%%%%%%%%%%%%%%%%%%%%%%%%%%%%%%%%%%%%%%%%%%%%%%%%%%%
%%%% KEEP THESE UPDATED ALSO IN THE MYMACROS.STY FILE   %%%%
%%%%%%%%%%%%%%%%%%%%%%%%%%%%%%%%%%%%%%%%%%%%%%%%%%%%%%%%%%%%

% Adding some operator names missing by default
\DeclareMathOperator{\cosec}{cosec}
\DeclareMathOperator{\powop}{pow} %\newcommand{\pow}{\operatorname{pow}}
\definecolor{frog}{rgb}{1,0,0}
\newcommand{\dx}{\,\text{d}x}
\newcommand{\du}{\,\text{d}u}
\newcommand{\dy}{\,\text{d}y}
\newcommand{\dt}{\,\text{d}t}

\KOMAoption{captions}{tableheading}
\makeatletter
\@ifpackageloaded{float}{}{\usepackage{float}}
\floatstyle{plain}
\@ifundefined{c@chapter}{\newfloat{video}{h}{lovid}}{\newfloat{video}{h}{lovid}[chapter]}
\floatname{video}{Video}
\newcommand*\listofvideos{\listof{video}{List of Videos}}
\makeatother
\makeatletter
\@ifpackageloaded{tcolorbox}{}{\usepackage[skins,breakable]{tcolorbox}}
\@ifpackageloaded{fontawesome5}{}{\usepackage{fontawesome5}}
\definecolor{quarto-callout-color}{HTML}{909090}
\definecolor{quarto-callout-note-color}{HTML}{0758E5}
\definecolor{quarto-callout-important-color}{HTML}{CC1914}
\definecolor{quarto-callout-warning-color}{HTML}{EB9113}
\definecolor{quarto-callout-tip-color}{HTML}{00A047}
\definecolor{quarto-callout-caution-color}{HTML}{FC5300}
\definecolor{quarto-callout-color-frame}{HTML}{acacac}
\definecolor{quarto-callout-note-color-frame}{HTML}{4582ec}
\definecolor{quarto-callout-important-color-frame}{HTML}{d9534f}
\definecolor{quarto-callout-warning-color-frame}{HTML}{f0ad4e}
\definecolor{quarto-callout-tip-color-frame}{HTML}{02b875}
\definecolor{quarto-callout-caution-color-frame}{HTML}{fd7e14}
\makeatother
\makeatletter
\@ifpackageloaded{caption}{}{\usepackage{caption}}
\AtBeginDocument{%
\ifdefined\contentsname
  \renewcommand*\contentsname{Table of contents}
\else
  \newcommand\contentsname{Table of contents}
\fi
\ifdefined\listfigurename
  \renewcommand*\listfigurename{List of Figures}
\else
  \newcommand\listfigurename{List of Figures}
\fi
\ifdefined\listtablename
  \renewcommand*\listtablename{List of Tables}
\else
  \newcommand\listtablename{List of Tables}
\fi
\ifdefined\figurename
  \renewcommand*\figurename{Figure}
\else
  \newcommand\figurename{Figure}
\fi
\ifdefined\tablename
  \renewcommand*\tablename{Table}
\else
  \newcommand\tablename{Table}
\fi
}
\@ifpackageloaded{float}{}{\usepackage{float}}
\floatstyle{ruled}
\@ifundefined{c@chapter}{\newfloat{codelisting}{h}{lop}}{\newfloat{codelisting}{h}{lop}[chapter]}
\floatname{codelisting}{Listing}
\newcommand*\listoflistings{\listof{codelisting}{List of Listings}}
\makeatother
\makeatletter
\makeatother
\makeatletter
\@ifpackageloaded{caption}{}{\usepackage{caption}}
\@ifpackageloaded{subcaption}{}{\usepackage{subcaption}}
\makeatother
\makeatletter
\@ifpackageloaded{sidenotes}{}{\usepackage{sidenotes}}
\@ifpackageloaded{marginnote}{}{\usepackage{marginnote}}
\makeatother
\makeatletter
\@ifpackageloaded{tcolorbox}{}{\usepackage[many]{tcolorbox}}
\makeatother
%%%% ---foldboxy preamble ----- %%%%%

\definecolor{fbx-default-color1}{HTML}{c7c7d0}
\definecolor{fbx-default-color2}{HTML}{a3a3aa}

\definecolor{fbox-color1}{HTML}{c7c7d0}
\definecolor{fbox-color2}{HTML}{a3a3aa}

% arguments: #1 typelabelnummer: #2 titel: #3
\newenvironment{fbx}[3]{\begin{tcolorbox}[enhanced, breakable,%
attach boxed title to top*={xshift=1.4pt},
boxed title style={boxrule=0.0mm, fuzzy shadow={1pt}{-1pt}{0mm}{0.1mm}{gray}, arc=.3em, rounded corners=east, sharp corners=west}, colframe=#1-color2, colbacktitle=#1-color1, colback = white, coltitle=black,  titlerule=0mm, toprule=0pt, bottomrule=.7pt, leftrule=.3em, rightrule=0pt, outer arc=.3em,  arc=0pt,	 sharp corners = east, left=.5em, bottomtitle=1mm, toptitle=1mm,title=\textbf{#2}\hspace{0.5em}{#3}]}
{\end{tcolorbox}}

% boxed environment with right border
\newenvironment{fbxSimple}[3]{\begin{tcolorbox}[enhanced, breakable,%
attach boxed title to top*={xshift=1.4pt},
boxed title style={boxrule=0.0mm, fuzzy shadow={1pt}{-1pt}{0mm}{0.1mm}{gray}, arc=.3em, rounded corners=east, sharp corners=west}, colframe=#1-color2, colbacktitle=#1-color1, colback = white, coltitle=black,  titlerule=0mm, toprule=0pt, bottomrule=.7pt, leftrule=.3em, rightrule=.7pt, outer arc=.3em,  	left=.5em, right=.5em, bottomtitle=1mm, toptitle=1mm,title=\textbf{#2}\hspace{0.5em}{#3}]}
{\end{tcolorbox}}

%%%% --- end foldboxy preamble ----- %%%%%
%%==== colors from yaml ===%
\definecolor{Example-color1}{HTML}{b2dac4}
\definecolor{Example-color2}{HTML}{004721}
\definecolor{Definition-color1}{HTML}{ffdc36}
\definecolor{Definition-color2}{HTML}{ffb948}
\definecolor{Theorem-color1}{HTML}{948bde}
\definecolor{Theorem-color2}{HTML}{584eab}
\definecolor{Task-color1}{HTML}{948bde}
\definecolor{Task-color2}{HTML}{584eab}
\definecolor{Weblink-color1}{HTML}{d18888}
\definecolor{Weblink-color2}{HTML}{660000}
\definecolor{Answer-color1}{HTML}{d18888}
\definecolor{Answer-color2}{HTML}{660000}
%=============%
\newcounter{quartocallouttipno}
\newcommand{\quartocallouttip}[1]{\refstepcounter{quartocallouttipno}\label{#1}}
\newcounter{quartocalloutcauno}
\newcommand{\quartocalloutcau}[1]{\refstepcounter{quartocalloutcauno}\label{#1}}
\usepackage{bookmark}
\IfFileExists{xurl.sty}{\usepackage{xurl}}{} % add URL line breaks if available
\urlstyle{same}
\hypersetup{
  colorlinks=true,
  linkcolor={blue},
  filecolor={Maroon},
  citecolor={Blue},
  urlcolor={Blue},
  pdfcreator={LaTeX via pandoc}}


\author{}
\date{}
\begin{document}


\chapter{Chapter 1 - Headings}\label{sec-headings}

Each row that begins with a \# is a chapter heading. So if each chapter
is a single file it's best to only use a single \# line at the top of
each document.

\section{Models for binary/binomial response}\label{sec-binary-response}

This chapter contains sections and subsections. This is the intro. Since
the chapter title begins with a \texttt{\#}, sections begin with
\texttt{\#\#}, and subsections should begin with \texttt{\#\#\#}.

Typical default numbering will look like this:

\begin{verbatim}
Chapter 1
Section 1.1
Subsection 1.1.1
Subsection 1.1.2
Section 1.2
Subsection 1.2.1
Subsection 1.2.2
Subsection 1.2.3
Section 1.3
\end{verbatim}

The \texttt{crossref} option in your global \texttt{.yml} file does
offer the chance to not number any chapters and just have sections in
each chapter, if you really don't want chapters.

\subsection{First subsection}\label{first-subsection}

Text goes here.

\subsection{Second subsection}\label{second-subsection}

Text goes here. And we can reference whole sections if we wish, like
Section~\ref{sec-binary-response}. You cannot link to other chapters as
they are standalones.

Let

\[D_3(y) \equiv \diff[3]{y}{x}\]

That was just an equation inside \texttt{\$\$} and \texttt{\$\$} tags.
For labelled equations you need to do the slightly unusual LaTeX
behaviour of using \texttt{\$\$} and then
\texttt{\textbackslash{}begin\{aligned\}} or similar inside. So unlike
in LaTeX you always need \texttt{\$\$} to start a math environment
(because the \texttt{\$\$} tells the browser to enable math mode). ie.
you need to use math-mode environments, not ones that themselves are
equivalent to beginning an equation.

\begin{equation}\phantomsection\label{eq-AROC}{
\begin{aligned}
\int_0^{\infty} x^3 \dx
\end{aligned}
}\end{equation}

This week we will learn how to model outcomes of interest that take one
of two categorical values (e.g.~yes/no, success/failure, alive/dead).
The independent responses \(Y_i\) can either be

\begin{itemize}
\item
  \emph{binary} (ungrouped), taking the value 1 (say success, with
  probability \(p_i\)) or 0 (failure, with probability \(1-p_i\)) or
\item
  \emph{binomial} (grouped), where \(Y_i\) is the number of successes in
  a given number of trials \(n_i\), with the probability of success
  being \(p_i\) and the probability of failure being \(1-p_i\).
\end{itemize}

In both cases the distribution of the \(Y_i\) is binomial, but in the
first case it is Bin\((1,p_i)\) and in the second case it is
Bin\((n_i,p_i)\).

\section{Binomial response}\label{binomial-response}

We will begin with models for binomial responses and we will look at
exploratory plots of the data, different choices of link function and
hypothesis tests about terms in the model. Finally we will examine
measures of goodness of fit of the model. Let's begin with an example.

\begin{tcolorbox}[enhanced jigsaw, toprule=.15mm, toptitle=1mm, rightrule=.15mm, bottomtitle=1mm, arc=.35mm, opacitybacktitle=0.6, opacityback=0, title=\textcolor{quarto-callout-tip-color}{\faLightbulb}\hspace{0.5em}{Beetle mortality}, coltitle=black, bottomrule=.15mm, breakable, colbacktitle=quarto-callout-tip-color!10!white, colframe=quarto-callout-tip-color-frame, titlerule=0mm, leftrule=.75mm, left=2mm, colback=white]

This is a ``Tip Callout'' box called via

\begin{verbatim}
:::{.callout-tip}
### Banner title like Bettle Mortality

Content here. This gives coloured boxes in HTML and LaTeX.
:::
\end{verbatim}

Now for a default code \texttt{panelset} as Quarto calls it. For
allowing dual code presentation.

\subsection{R}

\begin{Shaded}
\begin{Highlighting}[]
\NormalTok{beetles }\OtherTok{\textless{}{-}} \FunctionTok{read.csv}\NormalTok{(}\FunctionTok{url}\NormalTok{(}\StringTok{"http://www.stats.gla.ac.uk/\textasciitilde{}tereza/rp/beetles.csv"}\NormalTok{))}
\NormalTok{beetles}
\end{Highlighting}
\end{Shaded}

\begin{verbatim}
    dose number killed
1 1.6907     59      6
2 1.7242     60     13
3 1.7552     62     18
4 1.7842     56     28
5 1.8113     63     52
6 1.8369     59     53
7 1.8610     62     61
8 1.8839     60     60
\end{verbatim}

\subsection{Python}

\begin{Shaded}
\begin{Highlighting}[]
\ImportTok{import}\NormalTok{ pandas }\ImportTok{as}\NormalTok{ pd}
\NormalTok{beetles }\OperatorTok{=}\NormalTok{ pd.read\_csv(}\StringTok{"../resources/data/beetles.csv"}\NormalTok{)}
\NormalTok{beetles}
\end{Highlighting}
\end{Shaded}

\begin{verbatim}
     dose  number  killed
0  1.6907      59       6
1  1.7242      60      13
2  1.7552      62      18
3  1.7842      56      28
4  1.8113      63      52
5  1.8369      59      53
6  1.8610      62      61
7  1.8839      60      60
\end{verbatim}

\end{tcolorbox}

\begin{tcolorbox}[enhanced jigsaw, toprule=.15mm, toptitle=1mm, rightrule=.15mm, bottomtitle=1mm, arc=.35mm, opacitybacktitle=0.6, opacityback=0, title=\textcolor{quarto-callout-warning-color}{\faExclamationTriangle}\hspace{0.5em}{Warning \ref*{tip-special}: Example}, coltitle=black, bottomrule=.15mm, breakable, colbacktitle=quarto-callout-warning-color!10!white, colframe=quarto-callout-warning-color-frame, titlerule=0mm, leftrule=.75mm, left=2mm, colback=white]

\quartocallouttip{tip-special} 

Some text

\end{tcolorbox}

This box above can be linked to in the HTML via Tip~\ref{tip-special}.

\begin{tcolorbox}[enhanced jigsaw, toprule=.15mm, toptitle=1mm, rightrule=.15mm, bottomtitle=1mm, arc=.35mm, opacitybacktitle=0.6, opacityback=0, title=\textcolor{quarto-callout-warning-color}{\faExclamationTriangle}\hspace{0.5em}{Example}, coltitle=black, bottomrule=.15mm, breakable, colbacktitle=quarto-callout-warning-color!10!white, colframe=quarto-callout-warning-color-frame, titlerule=0mm, leftrule=.75mm, left=2mm, colback=white]

Some text

\end{tcolorbox}

\begin{tcolorbox}[enhanced jigsaw, toprule=.15mm, toptitle=1mm, rightrule=.15mm, bottomtitle=1mm, arc=.35mm, opacitybacktitle=0.6, opacityback=0, title=\textcolor{quarto-callout-warning-color}{\faExclamationTriangle}\hspace{0.5em}{Example}, coltitle=black, bottomrule=.15mm, breakable, colbacktitle=quarto-callout-warning-color!10!white, colframe=quarto-callout-warning-color-frame, titlerule=0mm, leftrule=.75mm, left=2mm, colback=white]

Some text

\end{tcolorbox}

\begin{tcolorbox}[enhanced jigsaw, toprule=.15mm, toptitle=1mm, rightrule=.15mm, bottomtitle=1mm, arc=.35mm, opacitybacktitle=0.6, opacityback=0, title=\textcolor{quarto-callout-warning-color}{\faExclamationTriangle}\hspace{0.5em}{Warning \ref*{tip-special2}: Example}, coltitle=black, bottomrule=.15mm, breakable, colbacktitle=quarto-callout-warning-color!10!white, colframe=quarto-callout-warning-color-frame, titlerule=0mm, leftrule=.75mm, left=2mm, colback=white]

\quartocallouttip{tip-special2} 

Some text

\end{tcolorbox}

Only the first and last of these had a label. They were
Tip~\ref{tip-special} and Tip~\ref{tip-special2}. Notice, however, that
you will need to define your own custom blocks. Since a \texttt{warning}
callout, which is used for an ``Example'' will suddenly be called
``Warning 1.1'' since the default naming of a labelled (and thus
numbered) block would need to be overridden.

Also notice the massive confusion caused by me using the
\texttt{\#tip-name} style label on a \texttt{warning} callout! Make sure
to use the right prefix in the reference else it'll do this.

\begin{tcolorbox}[enhanced jigsaw, toprule=.15mm, toptitle=1mm, rightrule=.15mm, bottomtitle=1mm, arc=.35mm, opacitybacktitle=0.6, opacityback=0, title=\textcolor{quarto-callout-caution-color}{\faFire}\hspace{0.5em}{Caution \ref*{cau-correct}: this is a caution callout}, coltitle=black, bottomrule=.15mm, breakable, colbacktitle=quarto-callout-caution-color!10!white, colframe=quarto-callout-caution-color-frame, titlerule=0mm, leftrule=.75mm, left=2mm, colback=white]

\quartocalloutcau{cau-correct} 

The reference name begins with \texttt{\#cau-}.

\end{tcolorbox}

Then when I reference it like this: Caution~\ref{cau-correct}, all is
well.




\end{document}
